\chapter{Reference tables}\label{app:tables}
This appendix collects, in the style of the Raven manual, every parameter, command, option, algorithm and output of the
groundwater coupling in one place. \emph{Default} is the value used when the item is not given; \emph{checked} is the range
Raven enforces at start-up (values outside it stop the run with a message); \emph{typical} ranges are guidance from the
literature \cite{freeze1979groundwater,domenico1998physical,johnson1967specific,calver2001riverbed}, not limits.

\begin{plain}
The tables here repeat, in one place, every parameter, command, option and output with its unit, default and allowed range
--- the same information as the chapters, arranged for looking up rather than reading.
\end{plain}

\section{Parameter classes}\label{app:classes}
As in Raven, every parameter belongs to a class, and the class decides where it is set and which part of the model it applies to.
\begin{longtable}{P{0.19\linewidth}P{0.13\linewidth}P{0.33\linewidth}P{0.25\linewidth}}
\caption{Parameter classes.}\\
\toprule\hdr Class & File & Applies to & Examples\\\midrule\endfirsthead
\caption[]{Parameter classes (continued)}\\
\toprule\hdr Class & File & Applies to & Examples\\\midrule\endhead
Aquifer class & \cmd{.rvp} & every layer made of that material & \cmd{K\_HORIZ}, \cmd{SPEC\_YIELD}\\
Aquifer profile & \cmd{.rvp} & every column whose dominant HRU has the profile; its layering & \cmd{SEEPAGE\_LEAKANCE}, \cmd{RIVERBED\_K}, layers\\
HRU & \cmd{.rvh} & one HRU & \cmd{AQUIFER\_PROFILE}, \cmd{AREA}, \cmd{ELEVATION}\\
Subbasin, channel & \cmd{.rvh}, \cmd{.rvp} & a subbasin's reach & channel profile, \cmd{Q\_REFERENCE}\\
Reservoir & \cmd{.rvh} & one reservoir & \cmd{:SeepageParameters}, \cmd{:AbsoluteCrest\allowbreak Height}\\
Groundwater option & \cmd{.rvg} & the whole groundwater model & \cmd{:GridCellSize}, \cmd{:SolverHead\allowbreak Tolerance}\\
Stress or observation & \cmd{.rvg}, \cmd{.rvt} & a well, boundary or monitoring well & \cmd{:Wells}, \cmd{:WellFlowRate}, \cmd{:BoundaryHead}\\
Initial condition & \cmd{.rvc} & the start of a run & \cmd{:InitialGWHeads}, hotstart blocks\\
Raven state variable & (computed) & an HRU store or flux & \cmd{GROUNDWATER}, \cmd{AET}, \cmd{DEPRESSION}\\\bottomrule
\end{longtable}

\section{Aquifer class parameters (\texttt{:AquiferClasses}, \texttt{.rvp})}\label{app:class}
\begin{longtable}{P{0.17\linewidth}P{0.23\linewidth}P{0.07\linewidth}P{0.09\linewidth}P{0.11\linewidth}P{0.19\linewidth}}
\caption{Aquifer class parameters (\cmd{:AquiferClasses}, \cmd{.rvp}).}\\
\toprule\hdr Parameter & Meaning & Units & Default & Checked & Typical range\\\midrule\endfirsthead
\caption[]{Aquifer class parameters (\cmd{:AquiferClasses}, \cmd{.rvp}) (continued)}\\
\toprule\hdr Parameter & Meaning & Units & Default & Checked & Typical range\\\midrule\endhead
\cmd{K\_HORIZ} & horizontal hydraulic conductivity & m/d & required & $>0$ & $10^{-7}$ (clay) -- $10^{3}$ (gravel)\\
\cmd{K\_VERT} & vertical hydraulic conductivity & m/d & $K_h/10$ & $>0$ & $K_h/100$ -- $K_h$\\
\cmd{SPEC\_STORAGE} & specific storage $S_s$ & 1/m & $10^{-5}$ & $\ge0$ & $10^{-6}$ (rock) -- $10^{-3}$ (soft clay)\\
\cmd{SPEC\_YIELD} & specific yield $S_y$ & -- & 0.1 & $[0,1]$ & 0.01 (clay) -- 0.35 (gravel)\\
\cmd{POROSITY} & saturated water content $\theta_s$ of the unsaturated columns (\cmd{:SoilZoneVadose}); not used by MODFLOW's flow & -- & 0.3 & $[S_y,1]$ & 0.01 (rock) -- 0.6 (clay)\\
\cmd{CLAPP\_B} & Clapp--Hornberger exponent $b$ of the retention curve \cite{clapp1978} (unsaturated columns) & -- & 4 & $>0$ & 4 (sand) -- 11 (clay)\\
\cmd{AIR\_ENTRY\_}\allowbreak\cmd{PRESSURE} & air-entry suction $\psi_a$ (unsaturated columns) & m & 0.3 & $>0$ & 0.05 (sand) -- 0.5 (clay)\\\bottomrule
\end{longtable}

\section{Typical values by material}\label{app:materials}
\begin{longtable}{P{0.207\linewidth}P{0.17\linewidth}P{0.15\linewidth}P{0.2\linewidth}P{0.16\linewidth}}
\caption{Typical values by material.}\\
\toprule\hdr Material & $K$ [m/d] & $S_y$ [--] & $S_s$ [1/m] & Porosity [--]\\\midrule\endfirsthead
\caption[]{Typical values by material (continued)}\\
\toprule\hdr Material & $K$ [m/d] & $S_y$ [--] & $S_s$ [1/m] & Porosity [--]\\\midrule\endhead
Gravel & $10^{1}$--$10^{3}$ & 0.13--0.35 & $10^{-6}$--$10^{-5}$ & 0.25--0.40\\
Sand, clean & $10^{0}$--$10^{2}$ & 0.10--0.35 & $10^{-5}$--$10^{-4}$ & 0.25--0.50\\
Sand, silty & $10^{-2}$--$10^{0}$ & 0.08--0.25 & $10^{-5}$--$10^{-4}$ & 0.30--0.45\\
Silt, loess & $10^{-3}$--$10^{0}$ & 0.03--0.20 & $10^{-4}$ & 0.35--0.50\\
Glacial till & $10^{-6}$--$10^{-1}$ & 0.02--0.15 & $10^{-4}$--$10^{-3}$ & 0.10--0.30\\
Clay & $10^{-7}$--$10^{-3}$ & 0.00--0.05 & $10^{-4}$--$10^{-2}$ & 0.40--0.70\\
Sandstone & $10^{-4}$--$10^{0}$ & 0.02--0.30 & $10^{-6}$--$10^{-5}$ & 0.05--0.30\\
Limestone, karst & $10^{-1}$--$10^{3}$ & 0.01--0.15 & $10^{-6}$ & 0.05--0.50\\
Fractured crystalline rock & $10^{-3}$--$10^{1}$ & 0.001--0.05 & $10^{-6}$ & 0.00--0.10\\\bottomrule
\end{longtable}

\section{Aquifer profile parameters (\texttt{:AquiferProfileParameters}, \texttt{.rvp})}\label{app:profile}
\begin{longtable}{P{0.28\linewidth}P{0.23\linewidth}P{0.06\linewidth}P{0.08\linewidth}P{0.09\linewidth}P{0.12\linewidth}}
\caption{Aquifer profile parameters (\cmd{:AquiferProfile\allowbreak Parameters}, \cmd{.rvp}).}\\
\toprule\hdr Parameter & Meaning & Units & Default & Checked & Typical range\\\midrule\endfirsthead
\caption[]{Aquifer profile parameters (\cmd{:AquiferProfile\allowbreak Parameters}, \cmd{.rvp}) (continued)}\\
\toprule\hdr Parameter & Meaning & Units & Default & Checked & Typical range\\\midrule\endhead
\cmd{INITIAL\_HEAD\_DEPTH} & initial depth to water below land & m & 5 & any & site data\\
\cmd{DEFAULT\_BEDROCK\_DEPTH} & bedrock depth where no bedrock map is given & m & 50 & $>0$ & 5--300\\
\cmd{SOIL\_ZONE\_DEPTH} & model top below land surface (ignored under \cmd{:SoilZone PHYSICAL}) & m & 0 & $\ge0$ & 0--3\\
\cmd{SEEPAGE\_LEAKANCE} & seepage conductance per m$^2$ of cell (not used for soil-zone HRUs) & 1/d & 1 & $\ge0$ & 0.1--10\\
\cmd{SEEPAGE\_SMOOTHING\_DEPTH} & depth below the top where seepage starts & m & 0.5 & $\ge0$ & 0.1--1\\
\cmd{RIVERBED\_K} & riverbed hydraulic conductivity & m/d & 0.5 & $\ge0$ & $10^{-3}$--10\\
\cmd{RIVERBED\_THICKNESS} & riverbed thickness & m & 1 & $>0$ & 0.3--3\\
\cmd{EXTINCTION\_DEPTH} & water-table ET extinction depth (0 = off) & m & 0 & $\ge0$ & 0.5--5 (root depth)\\
\cmd{RECHARGE\_DELAY} & recharge delay time constant (0 = off) & d & 0 & $\ge0$ & 0--365\\\bottomrule
\end{longtable}

\section{Layer types (\texttt{:AquiferProfiles}, \texttt{.rvp})}\label{app:layers}
\begin{longtable}{P{0.27\linewidth}P{0.3\linewidth}P{0.358\linewidth}}
\caption{Layer types (\cmd{:AquiferProfiles}, \cmd{.rvp}).}\\
\toprule\hdr Keyword & MODFLOW 6 setting & Behaviour\\\midrule\endfirsthead
\caption[]{Layer types (\cmd{:AquiferProfiles}, \cmd{.rvp}) (continued)}\\
\toprule\hdr Keyword & MODFLOW 6 setting & Behaviour\\\midrule\endhead
\cmd{AQUIFER}, \cmd{UNCONFINED\_AQUIFER} & \cmd{ICELLTYPE} 1, \cmd{ICONVERT} 1 & convertible: $S_y$ when unconfined, $S_s$ when full; transmissivity from saturated thickness\\
\cmd{CONFINED\_AQUIFER} & \cmd{ICELLTYPE} 0, \cmd{ICONVERT} 0 & always full: constant transmissivity, storage by $S_s$\\
\cmd{AQUITARD} & as \cmd{AQUIFER} & leaky layer with the (low) $K$ of its class\\
\cmd{AQUICLUDE} & \cmd{IDOMAIN} 0 & no flow; not allowed as top layer\\
thickness \cmd{TO\_BEDROCK} & layer bottom = bedrock & layers below it become pass-through cells\\
(automatic) pass-through & \cmd{IDOMAIN} $-1$ & layer thinner than 0.5 m or missing in the profile: connects the cells above and below\\\bottomrule
\end{longtable}

\section{Commands}\label{app:commands}
\begin{longtable}{P{0.08\linewidth}P{0.4\linewidth}P{0.448\linewidth}}
\caption{Commands of the groundwater coupling, by input file.}\\
\toprule\hdr
File & Command & Meaning (default)\\\midrule\endfirsthead
\caption[]{Commands of the groundwater coupling, by input file (continued).}\\
\toprule\hdr
File & Command & Meaning (default)\\\midrule\endhead
\cmd{.rvi} & \cmd{:GroundwaterModel NONE|MODFLOW6} & switch the coupling on (\cmd{NONE})\\
\cmd{.rvi} & \cmd{:rvg\_Filename file} & groundwater file, relative to the \cmd{.rvi} (\cmd{<model>.rvg})\\
\cmd{.rvh} & \cmd{AQUIFER\_PROFILE} column of \cmd{:HRUs} & profile name or \cmd{[NONE]}\\
\cmd{.rvh} & \cmd{:SeepageParameters} $k_{\mathrm{seep}}$ $h_{\mathrm{gw}}$ in \cmd{:Reservoir} & reservoir seepage constant; $h_{\mathrm{gw}}$ is replaced by the aquifer head under \cmd{:ReservoirExchange}\\
\cmd{.rvp} & \cmd{:AquiferClasses} \dots\ \cmd{:EndAquiferClasses} & material table (Section~\ref{app:class})\\
\cmd{.rvp} & \cmd{:AquiferProfiles} \dots\ \cmd{:EndAquiferProfiles} & name, number of layers, \{class, thickness or \cmd{TO\_BEDROCK}, type\}\\
\cmd{.rvp} & \cmd{:AquiferProfile\allowbreak Parameters} \dots & per-profile parameters (Section~\ref{app:profile})\\
\cmd{.rvg} & \cmd{:HRUGeometry file \{field\}} & HRU polygons, GeoJSON FeatureCollection, lon/lat or \cmd{:InputCRS} (\cmd{HRU\_ID}); required\\
\cmd{.rvg} & \cmd{:RiverGeometry file \{field\}} & river lines; subbasin field optional\\
\cmd{.rvg} & \cmd{:LandSurface file} & DEM, ESRI ASCII grid in lon/lat or \cmd{:InputCRS} (HRU elevations)\\
\cmd{.rvg} & \cmd{:InputCRS EPSG:code | TMERC | ALBERS | LCC \dots} & projection of every geometry and raster input file and of wells given by \cmd{X}, \cmd{Y} (lon/lat); Section~\ref{sec:inputcrs}\\
\cmd{.rvg} & \cmd{:BedrockSurface file} & bedrock elevation grid (\cmd{DEFAULT\_BEDROCK\_DEPTH})\\
\cmd{.rvg} & \cmd{:GridCellSize m} & cell size, $>0$ (1000)\\
\cmd{.rvg} & \cmd{:MinCellCoverage f} & coverage for an active column, $(0,1]$ (0.25)\\
\cmd{.rvg} & \cmd{:GridType REGULAR|QUADTREE|HRU\_MESH|FILE} & kind of grid (\cmd{REGULAR})\\
\cmd{.rvg} & \cmd{:MF6Simulation file}, \cmd{:MF6Model name} & existing model; which GWF model\\
\cmd{.rvg} & \cmd{:MF6StartDate yyyy-mm-dd} & date of the model's time zero (required)\\
\cmd{.rvg} & \cmd{:MF6CRS EPSG:code | TMERC | ALBERS | LCC \dots} & projection of the model's coordinates\\
\cmd{.rvg} & \cmd{:MF6GridFile file \{ID\}}, \cmd{:OverlapWeights} & cell polygons; or HRU--cell weights\\
\cmd{.rvg} & \cmd{:RavenTakesOver p \dots}, \cmd{:KeepPackages p \dots} & declared recharge/ET/UZF packages\\
\cmd{.rvg} & \cmd{:StreamPackage p \{AUX a | DOMINANT\_HRU\}} & RIV/DRN/GHB package to Raven's reaches\\
\cmd{.rvg} & \cmd{:SeepageToRaven ON|OFF} & \cmd{DRN\_RAVEN} at the model top (\cmd{ON})\\
\cmd{.rvg} & \cmd{:GridRotation} $\theta$ & degrees counterclockwise, $|\theta|\le360$; regular and quadtree grids (0)\\
\cmd{.rvg} & \cmd{:GridRefinement RIVERS|WELLS|LAKES s} or \cmd{LINES|POLYGONS file s} & target cell size $0<s<$ \cmd{:GridCellSize}; not for \cmd{FILE}; \cmd{LAKES}: every reservoir's lake HRU\\
\cmd{.rvg} & \cmd{:GridFile file \{ID\}} & cell polygons for \cmd{:GridType FILE} (\cmd{CELL\_ID})\\
\cmd{.rvg} & \cmd{:QuadtreeMaxLevel n} & $0\le n\le12$; \cmd{QUADTREE} only (from refinement sizes)\\
\cmd{.rvg} & \cmd{:FlowCorrection AUTO|XT3D|XT3D\_RHS|NONE} & XT3D (\cmd{AUTO}: off)\\
\cmd{.rvg} & \cmd{:LayerConnection INDEX|ELEVATION} & layers connected by number or by elevation (\cmd{INDEX})\\
\cmd{.rvg} & \cmd{:MF6Library path} & MODFLOW 6 library (else \cmd{RAVEN\_MF6\_LIB}, next to the Raven executable, then \cmd{lib/mf6/} of the Raven checkout, then the system path)\\
\cmd{.rvg} & \cmd{:SeepageReturnTo SV} & store receiving seepage (lowest soil layer; \cmd{SURFACE\_WATER} for soil-less HRUs)\\
\cmd{.rvg} & \cmd{:RiverBedDepth m|REFERENCE\_FLOW} & riverbed depth below the DEM (\cmd{REFERENCE\_FLOW})\\
\cmd{.rvg} & \cmd{:GWInitialization STEADY\_STATE r} & steady-state start, recharge $r\ge0$ mm/d (off); if it does not converge, the initial heads are used\\
\cmd{.rvg} & \cmd{:SurfaceWater\allowbreak Exchange SV group LEAKANCE L} & store--aquifer exchange, leakance $\Lambda>0$ 1/d (generated and existing models)\\
\cmd{.rvg} & \cmd{:SoilZone CONCEPTUAL | PHYSICAL SOIL[0] \dots\ SOIL[$M-1$]} & soil zone owned by Raven, water table may rise into it (Section~\ref{ch:soilzone}); the last \cmd{:SoilZone} counts (\cmd{CONCEPTUAL})\\
\cmd{.rvg} & \cmd{:SoilZoneIteration n \{$\delta_h$\} \{$\delta_q$\}} & interface updates after MODFLOW converged, $n\ge0$ (0: explicit), head [m] and flow [--] tolerances (100, 0.0001, 0.000001)\\
\cmd{.rvg} & \cmd{:SoilZoneSubgrid n} & DEM points per cell side for the saturated area, $1\le n\le20$ (5, i.e.\ 25 points per cell)\\
\cmd{.rvg} & \cmd{:SoilZoneOutput group} & \cmd{GWSoilExchange.csv} only for this HRU group (all soil-zone HRUs)\\
\cmd{.rvg} & \cmd{:SoilZoneVadose NONE | RICHARDS \{$\epsilon$\} \{N\} | KINEMATIC \{$\epsilon$\} \{N\} | LAG \{$t_m$\} \{N\}} & unsaturated columns below the soil zone: Richards' equation (Section~\ref{sec:szvadose}), kinematic wave (Section~\ref{sec:szvadosekw}) or linear storage cascade (Section~\ref{sec:szvadoselag}); largest content change in a substep $0<\epsilon\le0.2$, travel time per metre $t_m\ge0$ [d/m] (0: derived per material), nodes $3\le N\le200$ (NONE; $\epsilon$ 0.02, $t_m$ 0, $N$ 16)\\
\cmd{.rvg} & \cmd{:SoilZoneVadoseClasses K} & unsaturated columns per exchange column, one per depth class of its DEM samples (Section~\ref{sec:vzclasses}); $1\le K\le10$ (1); with \cmd{:SoilZoneVadose NONE} reported as ignored\\
\cmd{.rvg} & \cmd{:SoilZoneVadoseInitialFlux q} & steady start of the unsaturated columns for the recharge $q\ge0$ [mm/d] on a cold start: each column at its steady profile for the flux it receives, $(1-F)\,q$ ($F$ of its class at the initial head); a column that cannot carry it at its backed-up steady state (warning); ignored with a warning on a groundwater hotstart and with \cmd{:SoilZoneVadose NONE} (0: at rest)\\
\cmd{.rvg} & \cmd{:SoilZoneCapillaryRise ON | OFF} & capillary rise into the soil zone (Section~\ref{sec:szcaprise}) (ON with an unsaturated column: \cmd{RICHARDS}, \cmd{KINEMATIC} or \cmd{LAG}; OFF with \cmd{NONE})\\
\cmd{.rvg} & \cmd{:SoilZoneExchangeTiming LAGGED | AFTER} & order of Raven's processes and the exchange (Section~\ref{sec:sztiming}) (\cmd{LAGGED})\\
\cmd{.rvg} & \cmd{:SoilZoneClasses K} & Raven's processes per class of columns, $1\le K\le50$ (Section~\ref{sec:szclasses}) (1)\\
\cmd{.rvg} & \cmd{:SoilZoneAquiferClass class} & existing models: \cmd{CLAPP\_B}, \cmd{AIR\_ENTRY\_PRESSURE}, \cmd{POROSITY} of the cells for the unsaturated columns (4, 0.3 m, $S_y+0.1$), and \cmd{SPEC\_YIELD} for a steady model without one (0.2); reported as ignored on a model Raven builds\\
\cmd{.rvg} & \cmd{:ReservoirExchange} & reservoir--aquifer exchange, generated and existing models (off)\\
\cmd{.rvg} & \cmd{:Wells} \dots\ \cmd{:EndWells} & ID, \cmd{NAME, LATITUDE, LONGITUDE, SCREEN\_TOP, SCREEN\_BOT} (\cmd{X, Y} instead of \cmd{LATITUDE, LONGITUDE} with \cmd{:InputCRS})\\
\cmd{.rvg} & \cmd{:ObservationWells} \dots & ID, \cmd{NAME, LATITUDE, LONGITUDE, SCREEN\_ELEV} (or \cmd{X, Y}, as for \cmd{:Wells})\\
\cmd{.rvg} & \cmd{:GeneralHeadBoundary n file HEAD h CONDUCTANCE c \{LAYERS ALL|TOP\}} & head-dependent boundary; $c>0$ m$^2$/d per cell (\cmd{ALL})\\
\cmd{.rvg} & \cmd{:SpecifiedHeadBoundary n file HEAD h \{LAYERS ALL|TOP\}} & fixed-head boundary (\cmd{ALL})\\
\cmd{.rvg} & \cmd{:SolverHead\allowbreak Tolerance m} & Newton target tolerance, $>0$ (0.00001)\\
\cmd{.rvg} & \cmd{:SolverMax\allowbreak OuterIterations n} & Newton iteration limit, $\ge1$ (500)\\
\cmd{.rvg} & \cmd{:HeadSaveFrequency n} & steps between saved head fields, $\ge1$ (30)\\
\cmd{.rvg} & \cmd{:WriteNetCDFHeads} & write \cmd{GWHeads.nc} (off; needs a NetCDF build)\\
\cmd{.rvg} & \cmd{:WriteGWBudgetFile} & MODFLOW budget file \cmd{mf6/gwf.cbc} every \cmd{:HeadSaveFrequency} steps, with specific discharge and saturation, and \cmd{mf6/gwf.porosity}, for particle tracking (Section~\ref{sec:budgetfile}); generated models only (off)\\
\cmd{.rvg} & \cmd{:LinkageCache file} & linkage cache (\cmd{<output>/mf6/linkage.cache})\\
\cmd{.rvt} & \cmd{:WellFlowRate id \{units\}} \dots\ \cmd{:EndWellFlowRate} & well rate series, negative = pumping; units \cmd{m3/d} (default), \cmd{m3/s} or \cmd{L/s} (converted to m$^3$/d); others stop the run\\
\cmd{.rvt} & \cmd{:BoundaryHead name \{units\}} \dots & time-varying boundary head; units \cmd{m} (default); others stop the run\\
\cmd{.rvt} & \cmd{:ObservationData GW\_HEAD id units} \dots & observed heads for diagnostics\\
\cmd{.rvc} & \cmd{:InitialGWHeads FROM\_PROFILE|DEPTH\_BELOW\_SURFACE d|ELEVATION z} & initial heads (\cmd{FROM\_PROFILE})\\
\cmd{.rvc} & \cmd{:GWHeads}, \cmd{:GWRechargeStore}, \cmd{:GWReservoirSeepage}, \cmd{:GWRiverLossCarry}, \cmd{:GWSoilZoneColumns}, \cmd{:GWVadoseColumns}, \cmd{:GWVadoseClassColumns}, \cmd{:GWSoilZoneLag} & hotstart blocks written by Raven (Section~\ref{app:hotstart})\\\bottomrule
\end{longtable}

\section{Options}\label{app:options}
\begin{longtable}{P{0.28\linewidth}P{0.28\linewidth}P{0.368\linewidth}}
\caption{Options of the groundwater commands.}\\
\toprule\hdr Command & Option & Effect\\\midrule\endfirsthead
\caption[]{Options of the groundwater commands (continued)}\\
\toprule\hdr Command & Option & Effect\\\midrule\endhead
\cmd{:GroundwaterModel} & \cmd{NONE} & Raven without groundwater: results identical to Raven 4.15 except for models with reservoirs, ensembles and age-tracer hotstarts (Table~\ref{tab:ravendiff})\\
 & \cmd{MODFLOW6} & coupled HRUs exchange water with MODFLOW 6\\
\cmd{:InitialGWHeads} & \cmd{FROM\_PROFILE} & $z_{\mathrm{land}}-$\cmd{INITIAL\_HEAD\_DEPTH} of the column's profile\\
 & \cmd{DEPTH\_BELOW\_SURFACE d} & $z_{\mathrm{land}}-d$ everywhere\\
 & \cmd{ELEVATION z} & $z$ everywhere (at least 0.1 m above the lowest active bottom)\\
\cmd{:GWInitialization} & \cmd{STEADY\_STATE r} & steady state with recharge $r$ over covered areas before day 1; skipped on hotstart\\
\cmd{:RiverBedDepth} & value & bed = lowest DEM along the river $-$ value\\
 & \cmd{REFERENCE\_FLOW} & bed = lowest DEM $-$ Raven channel depth at the reference flow (max 20 m)\\
\cmd{LAYERS} (boundaries) & \cmd{ALL} & every active layer whose bottom is below the head\\
 & \cmd{TOP} & the first such layer only\\
\cmd{:SeepageReturnTo} & any storage state variable & e.g.\ \cmd{SOIL[2]}, \cmd{SURFACE\_WATER}, \cmd{DEPRESSION}\\\bottomrule
\end{longtable}

\section{Exchange algorithms}\label{app:algorithms}
\begin{longtable}{P{0.13\linewidth}P{0.17\linewidth}P{0.35\linewidth}P{0.26\linewidth}}
\caption{Exchange algorithms and their safeguards.}\\
\toprule\hdr Exchange & MODFLOW package (name) & Formulation & Safeguard\\\midrule\endfirsthead
\caption[]{Exchange algorithms and their safeguards (continued)}\\
\toprule\hdr Exchange & MODFLOW package (name) & Formulation & Safeguard\\\midrule\endhead
Recharge & RCH (\cmd{RCH\_RAVEN}) & $V_k=10^{6}A_k\,r_k/1000$; optional linear delay; spread by overlap & fixed-head cells skipped; exact per HRU\\
Seepage & DRN (\cmd{DRN\_RAVEN}) & $C f(h)(h-z_{\mathrm{top}}+d_s)$, $C=\Lambda_s|\Omega_i|$ & returned by overlap; soil-less HRUs to \cmd{SURFACE\_WATER}\\
Rivers & RIV (\cmd{RIV\_RAVEN}) & $C(s-h)$, or $C(s-z_{\mathrm{bed}})$ below the bed; $C=K_bW_pL_{p,i}/M_b$ & $C$ limited by channel flow; unsupplied loss carried\\
Water-table ET & EVT (\cmd{EVT\_RAVEN}) & unmet PET over the cell, linear to $d_x$ & only ET-enabled HRUs; never above unmet demand\\
Wetlands, lakes & RIV (\cmd{SWX\_RAVEN}) & bed = top, stage = top + store depth & $C$ limited by stored water; one exchange per store\\
Reservoirs & WEL (\cmd{RES\_RAVEN}) & Raven seepage law, $h_{\mathrm{gw}}=\sum w_i\max(z^{\mathrm{wt}}_i,z_{\mathrm{top},i})$; one-step lag & in-transit volume reported and saved\\
Wells & WEL (\cmd{WEL\_RAVEN}) & rate split by transmissivity of screened layers & \cmd{AUTO\_FLOW\_REDUCE}\\
Head boundary & GHB (\cmd{GHB\_RAVEN}) & $C_g(h_g-h)$ & heads kept above cell bottoms\\
Fixed head & CHD (\cmd{CHD\_RAVEN}) & head held at $h_c$ & heads kept above cell bottoms\\\bottomrule
\end{longtable}

\section{Solver settings (\texttt{sim.ims})}\label{app:solver}
\begin{longtable}{P{0.34\linewidth}P{0.2\linewidth}P{0.388\linewidth}}
\caption{Solver settings written to \cmd{sim.ims}.}\\
\toprule\hdr Setting & Value & Note\\\midrule\endfirsthead
\caption[]{Solver settings written to \cmd{sim.ims} (continued)}\\
\toprule\hdr Setting & Value & Note\\\midrule\endhead
Formulation & Newton & in \cmd{gwf.nam}, without its \cmd{UNDER\_RELAXATION} keyword (Section~\ref{ch:mf6})\\
\cmd{OUTER\_DVCLOSE} & 0.00001 m & target; raised to 10$\times$ for the rest of a step after 100 iterations (or half the limit)\\
\cmd{OUTER\_MAXIMUM} & 500 & \cmd{:SolverMax\allowbreak OuterIterations}\\
\cmd{UNDER\_RELAXATION} & \cmd{DBD} & theta 0.7, kappa 0.07, gamma 0.1\\
\cmd{BACKTRACKING\_NUMBER} & 20 & tolerance 1.05, reduction 0.1, residual limit 0.002\\
\cmd{INNER\_MAXIMUM} & 500 & linear iterations\\
\cmd{INNER\_DVCLOSE}, \cmd{INNER\_RCLOSE} & $0.1\times$ outer, 1 m$^3$/d & \\
\cmd{LINEAR\_ACCELERATION} & \cmd{BICGSTAB} & ILU, 5 levels, drop tolerance $10^{-4}$\\\bottomrule
\end{longtable}

\section{Output: \texttt{GWBudget.csv} columns}\label{app:outputs}
\begin{longtable}{P{0.44\linewidth}P{0.5\linewidth}}
\caption{Columns of \cmd{GWBudget.csv}.}\\
\toprule\hdr
Column & Meaning\\\midrule\endfirsthead
\caption[]{Columns of \cmd{GWBudget.csv} (continued).}\\
\toprule\hdr
Column & Meaning\\\midrule\endhead
\cmd{time [d]}, \cmd{date} & end of the step (with time of day for sub-daily steps)\\
\cmd{Raven recharge [m3/d]} & recharge leaving Raven, before any delay\\
\cmd{MF6 recharge [m3/d]} & recharge entering MODFLOW\\
\cmd{seepage to Raven [m3/d]} & drain discharge at the model top\\
\cmd{storage release [m3/d]} & water released from aquifer storage (negative: stored)\\
\cmd{MF6 balance error [m3/d]}, \cmd{[\%]} & MODFLOW imbalance; \% of gross flow\\
\cmd{returned to Raven [m3]} & seepage added to Raven stores\\
\cmd{outer iterations}, \cmd{converged} & Newton iterations; 1 target, 2 relaxed, 0 not converged\\
\cmd{recharge delay storage [m3]} & water in delay reservoirs\\
\cmd{recharge onto CHD cells [m3/d]} & recharge of HRUs lying wholly on fixed-head cells\\
\cmd{river leakage to aquifer [m3/d]}, \cmd{aquifer discharge to rivers [m3/d]} & river exchange by direction\\
\cmd{river exchange applied to Raven [m3]} & net volume applied to Raven reaches\\
\cmd{river loss carried to next step [m3]} & unsupplied loss taken next step (normally 0)\\
\cmd{water-table ET [m3/d]} & evaporation from the water table\\
\cmd{wells specified [m3/d]}, \cmd{wells actual [m3/d]} & requested and delivered well rates\\
\cmd{GHB [m3/d]}, \cmd{CHD [m3/d]} & boundary flows (positive into the aquifer)\\
\cmd{surface-store leakage to aquifer [m3/d]}, \cmd{aquifer discharge to surface stores [m3/d]} & wetland/lake exchange\\
\cmd{surface-store exchange applied to Raven [m3]} & net volume applied to Raven stores\\
\cmd{reservoir seepage sent [m3/d]}, \cmd{delivered [m3/d]}, \cmd{awaiting delivery [m3]} & reservoir exchange\\
\cmd{reservoir gain not supplied by MODFLOW [m3]} & lake gain a drying cell could not supply; taken back from the reach (normally 0)\\
\cmd{other model packages [m3/d]} & existing models only: net flow of the model's own packages (wells, boundaries, \dots), positive into the aquifer\\
\cmd{soil zone to aquifer [m3/d]}, \cmd{aquifer to soil zone [m3/d]} & \cmd{:SoilZone PHYSICAL}: soil-zone exchange by direction\\
\cmd{rejected recharge [m3]}, \cmd{soil-zone overflow [m3]} & recharge kept in the soil on saturated ground; water a full soil zone sent to \cmd{SURFACE\_WATER}\\
\cmd{soil-zone exchange applied to Raven [m3]}, \cmd{soil-zone shortfall [m3]} & rejected recharge plus net exchange applied to the soil; water MODFLOW took that the soil did not hold (0)\\
\cmd{interface updates}, \cmd{interface converged}, \cmd{interface change [m]} & soil-zone updates in the step, 1 if the interface converged, last change of the soil level\\
\cmd{soil-zone layer redistribution [m3]} & percolation Raven computed on the HRU mean that no soil-zone column could take (it stays in the layer above)\\
\cmd{seepage and soil inflow withheld [m3]}, \cmd{river gains withheld [m3]}, \cmd{surface-store gains withheld [m3]} & discharge not handed to Raven for the soil-zone shortfall and, on unconverged steps, MODFLOW's negative residual (0)\\
\cmd{negative recharge the soil zone did not hold [m3]} & the cut of a negative recharge (Section~\ref{sec:negrecharge}) that the soil zone could not hold (0)\\
\cmd{unsaturated zone inflow [m3]}, \cmd{unsaturated zone to aquifer [m3]} & \cmd{:SoilZoneVadose} other than \cmd{NONE}: water entering the unsaturated columns; water they pass to the water table (including the excess a rising water table absorbs)\\
\cmd{capillary rise to soil zone [m3]}, \cmd{unsaturated zone storage [m3]} & (same) rise from the columns' top nodes into the soil zone (\cmd{RICHARDS}); water in the columns relative to their reference profile at the end of the step\\
\cmd{capillary rise from the water table [m3]} & capillary rise on: rise from the water table into the soil-zone layers (Section~\ref{sec:szcaprise})\\
\cmd{lagged-exchange debt [m3]}, \cmd{lagged-exchange repayment [m3]} & \cmd{LAGGED} timing: lent water Raven's processes used that the step's exchange did not supply; the aquifer's repayment in the step (Section~\ref{sec:sztiming})\\
\cmd{soil zone to specified heads [m3/d]}, \cmd{specified heads to soil zone [m3/d]} & soil-zone columns over specified heads, or an existing model (Section~\ref{sec:szchd})\\
\cmd{recharge not applied by MODFLOW [m3/d]} & an existing model without \cmd{NEWTON}: recharge MODFLOW did not apply (dry cells, specified heads), signed\\
\cmd{reservoir seepage not delivered by MODFLOW [m3]} & an existing model with \cmd{:ReservoirExchange}: reservoir seepage MODFLOW did not take (cells under the lake that are specified heads or inactive), returned to the reservoir's subbasin in the same step (normally 0)\\\bottomrule
\end{longtable}
Other output files: \cmd{GWSoilExchange.csv} (\cmd{:SoilZone PHYSICAL}, per soil-zone HRU and step: soil-zone base, soil
water level, water table, saturated fraction, exchange with the aquifer both ways, rejected recharge and overflow [mm], soil
water and capacity [mm], interface updates and change, and with specified heads the exchange with them both ways; the
exchange columns summed over the HRUs equal those of \cmd{GWBudget.csv}; Table~\ref{tab:outputs}), \cmd{GWHRUState.csv} (water table and depth per coupled HRU), \cmd{GWHeads.csv} (monitoring wells;
$-1.2345$ when dry), \cmd{GWConnectivity.csv} (wet areas, wet fraction per profile), \cmd{GWModelSummary.txt}, \cmd{GWHeads.nc}
(optional), \cmd{mf6/} (MODFLOW model, list files, binary heads).

\paragraph{Which columns are written.} The columns appear in the order of the table, and only where they apply:
\cmd{other model packages} with an existing model; the soil-zone columns, from \cmd{soil zone to aquifer} to \cmd{negative
recharge the soil zone did not hold}, with \cmd{:SoilZone PHYSICAL}; the unsaturated-zone columns with a
\cmd{:SoilZoneVadose} method; \cmd{capillary rise from the water table} with capillary rise on; the lagged-exchange columns
with \cmd{LAGGED} timing; the specified-head columns where soil-zone columns lie over specified heads; and, last,
\cmd{recharge not applied by MODFLOW} for an existing model without \cmd{NEWTON} (signed: negative where a negative recharge
was taken out of \cmd{GROUNDWATER} without MODFLOW supplying it), then \cmd{reservoir seepage not delivered by MODFLOW} for an
existing model with \cmd{:ReservoirExchange}. \cmd{MF6 balance error [\%]} divides by the gross flow
(Section~\ref{ch:balance}). Every row is stamped with the end of its step.

\paragraph{Run totals.} At the end of a run Raven appends to \cmd{GWModelSummary.txt}: time steps and MODFLOW-unconverged steps; with
soil-zone links the interface-unconverged steps, the columns' substeps (and Newton iterations except with \cmd{LAG}), the outer
iterations and steps in which drained columns were kept 1~mm above their bottoms; with an interface iteration the
supply-limited link updates and the largest take beyond the free water, and the connection-test drops (link updates, steps,
and those followed by a plain update); without \cmd{:SoilZone PHYSICAL} the discharge withheld on unconverged steps for
MODFLOW's negative residual (only when such steps occurred); the negative recharge cut and the part not held (only when not zero); for an existing
model with \cmd{:ReservoirExchange} the reservoir seepage not delivered, and for an existing model without \cmd{NEWTON} the recharge
not applied (both always written, also 0).


\section{Hotstart blocks (\texttt{solution.rvc})}\label{app:hotstart}
\begin{longtable}{P{0.3\linewidth}P{0.64\linewidth}}
\caption{Hotstart blocks written to \cmd{solution.rvc}.}\\
\toprule\hdr Block & Content\\\midrule\endfirsthead
\caption[]{Hotstart blocks written to \cmd{solution.rvc} (continued)}\\
\toprule\hdr Block & Content\\\midrule\endhead
\cmd{:GWHeads nlay nrow ncol} & MODFLOW head of every cell (raw values, 12 significant digits; 17 with \cmd{:SoilZone PHYSICAL}); checked against the grid (counts that are not positive, or whose product exceeds $10^9$, stop the run with a message)\\
\cmd{:GWRechargeStore} & HRU ID, water in the recharge-delay reservoir [m$^3$]\\
\cmd{:GWReservoirSeepage} & subbasin ID, reservoir seepage in transit [m$^3$]\\
\cmd{:GWRiverLossCarry} & subbasin ID, unsupplied river loss to take next step [m$^3$] (only when non-zero)\\
\cmd{:GWSoilZoneColumns} & \cmd{:SoilZone PHYSICAL}: HRU ID, cell (grid column number; 0 for the rest of the HRU), water of each soil-zone layer from the bottom up [mm]; a block that does not match the HRU's columns is reported and the columns start from the water table\\
\cmd{:GWVadoseColumns} & \cmd{:SoilZoneVadose RICHARDS}, \cmd{KINEMATIC} or \cmd{LAG} (interchangeable between the methods): HRU ID, cell (grid column number), water table [m], water contents of the column's nodes from the top [--]\\
\cmd{:GWVadoseClassColumns} & \cmd{:SoilZoneVadoseClasses} $K>1$: HRU ID, cell (grid column number), depth class, classes of the cell, water table [m], water contents of the class column's nodes from the top [--]\\
\cmd{:GWSoilZoneLag} & \cmd{LAGGED} timing: HRU ID, cell, exchange of the last step into the aquifer [m$^3$/d], debt the aquifer still repays [m$^3$]\\
\cmd{:AgeTracerDatum date time} & with an age tracer (Raven core, also without groundwater), when the datum of the tracer's entry times is before the run's start: that datum (the start of the first run of a chain of restarts); \cmd{:StartTime} is the start of the run that wrote the file\\\bottomrule
\multicolumn{2}{p{\dimexpr\linewidth-2\tabcolsep}}{\footnotesize Each block ends with the matching \cmd{:End} line (\cmd{:EndGWHeads}, \cmd{:EndGWRechargeStore}, \cmd{:EndGWReservoirSeepage}, \cmd{:EndGWRiverLossCarry}, \cmd{:EndGWSoilZoneColumns}, \cmd{:EndGWVadoseColumns}, \cmd{:EndGWVadoseClassColumns}, \cmd{:EndGWSoilZoneLag}). In a run without \cmd{:GroundwaterModel MODFLOW6} these blocks are ignored with a warning. A value that is not a finite number in any block, or negative water in \cmd{:GWSoilZoneColumns}, \cmd{:GWVadoseColumns} or \cmd{:GWVadoseClassColumns}, stops the run with a message naming the block.}\\
\end{longtable}

\section{Error messages}\label{app:errors}
All stop the run with the message shown (abbreviated) in \cmd{Raven\_errors.txt}.
\begin{longtable}{P{0.5\linewidth}P{0.44\linewidth}}
\caption{Error messages and remedies.}\\
\toprule\hdr
Message & Remedy\\\midrule\endfirsthead
\caption[]{Error messages and remedies (continued).}\\
\toprule\hdr
Message & Remedy\\\midrule\endhead
\cmd{:HRUGeometry is required in the .rvg file} & add the HRU polygon file\\
cannot open GeoJSON / raster file; missing ID field & check paths and the ID field name\\
HRU \dots\ has an AQUIFER\_PROFILE but no polygon & polygon IDs must match \cmd{.rvh} IDs\\
HRU \dots\ references unknown AQUIFER\_PROFILE & define the profile in the \cmd{.rvp}\\
HRU \dots\ does not overlap any active groundwater cell & smaller \cmd{:GridCellSize} or \cmd{:MinCellCoverage}; check bedrock\\
no HRU has an AQUIFER\_PROFILE; coupled model requires \cmd{:AquiferProfiles} & fill the \cmd{.rvh} column; define profiles\\
no process moves water into GROUNDWATER & route recharge to \cmd{GROUNDWATER}\\
process \# removes water from GROUNDWATER in HRU \dots & restrict the process with \cmd{:-->Conditional}\\
aquifer class/profile \dots\ must be positive / may not be negative / between 0 and 1 & correct the value (Sections~\ref{app:class}, \ref{app:profile})\\
top layer cannot be an AQUICLUDE; layer thickness must be positive & correct the profile\\
\cmd{:MinCellCoverage}, \cmd{:GridCellSize}, \cmd{:SolverHead\allowbreak Tolerance}, \cmd{:SolverMax\allowbreak OuterIterations}, \cmd{:HeadSaveFrequency} out of range & correct the value (Section~\ref{app:commands})\\
improper format of command \dots & correct the command syntax\\
command \dots\ in \dots: '\emph{token}' is not a number / is not a whole number / is not a whole-number ID & give a plain number, in the units of the manual, without units or thousands separators\\
command \dots\ has the surplus value '\emph{token}' & remove it; a comma separates values (2000.0, not 2,000.0)\\
boundary \dots: the values after the file must come in keyword-value pairs / unknown keyword & \cmd{HEAD}, \cmd{CONDUCTANCE} and \cmd{LAYERS}, each with its value\\
\cmd{:GridCellSize} \dots\ m gives about \dots\ cells per layer over the HRUs' extent \dots, more than 1e+07 / the \cmd{:GridRefinement} sizes give a regular grid of about \dots\ cells per layer & give the sizes in metres as plain numbers; for refinement around features only use \cmd{QUADTREE} or \cmd{HRU\_MESH}\\
\cmd{:WellFlowRate} \emph{id}: units \dots\ are not supported / \cmd{:BoundaryHead} \emph{name}: units \dots & rates in \cmd{m3/d}, \cmd{m3/s} or \cmd{L/s}; heads in \cmd{m}\\
boundary \dots\ needs a HEAD value / positive CONDUCTANCE & complete the boundary command\\
two wells / observation wells / boundaries have the same ID or name & make IDs unique\\
\cmd{:WellFlowRate} \emph{id} / \cmd{:BoundaryHead} \emph{name} in the .rvt file names no well / boundary & correct the ID or name (the message lists the wells or boundaries of the .rvg file)\\
\cmd{:MF6Model}, \cmd{:MF6StartDate}, \cmd{:KeepPackages}, \cmd{:SeepageToRaven} apply only to an existing model & remove them from a model Raven builds\\
unknown aquifer profile parameter \dots & correct the \cmd{:Attributes} name (Section~\ref{app:commands})\\
boundary \dots: LAYERS must be ALL or TOP & correct the keyword\\
HRU \dots\ is in two \cmd{:SurfaceWater\allowbreak Exchange} groups for the same store & one exchange per store and HRU\\
well \dots\ is outside the active groundwater model & move the well or enlarge the coupled area\\
\cmd{:GWHeads} does not match the groundwater grid & hotstart file from a different grid\\
the model's \emph{type} package \emph{name} would add recharge or ET beside Raven's & list it in \cmd{:RavenTakesOver} or \cmd{:KeepPackages}\\
the model name file \dots\ has a \emph{type} package named \emph{name}, a name Raven reserves for the packages it adds (\cmd{RCH\_RAVEN}, \cmd{DRN\_RAVEN}, \cmd{EVT\_RAVEN}, \cmd{SWX\_RAVEN}, \cmd{RES\_RAVEN}, \cmd{SZX/SZQ/SZD/SZW\_RAVEN}) & rename the package, or remove its line if Raven added it to a model it wrote; a recharge or ET package may instead be listed in \cmd{:RavenTakesOver}\\
\cmd{:MF6StartDate} \dots / the model ends before Raven's run & correct the dates or \cmd{:Duration}\\
stress period \dots\ is not a whole number of Raven time steps & choose a Raven time step that divides the periods\\
cells of the grid file differ from the model's cell area & wrong projection, length unit or cell numbering in the file\\
give one of \cmd{:MF6CRS}, \cmd{:MF6GridFile} or \cmd{:OverlapWeights} & link HRUs in exactly one way\\
the period data of \emph{package} cannot be renumbered safely & start Raven at \cmd{:MF6StartDate}\\
\emph{package} uses time-array series (TAS6) / time series, \dots\ and the simulation holds several models & start Raven at \cmd{:MF6StartDate}\\
TDIS period data given as INTERNAL / the water mover \dots\ gives its period data as INTERNAL & list the rows in the file itself or in an \cmd{OPEN/CLOSE} file\\
an aquifer profile has EXTINCTION\_DEPTH $>0$ \dots\ and the model keeps its EVT/UZF package & list that package in \cmd{:RavenTakesOver}, or set \cmd{EXTINCTION\_DEPTH} 0\\
with \cmd{:MF6Simulation} \dots\ remove \dots\ (these apply to models Raven builds) & remove the generated-model commands (grid options, \cmd{:LandSurface}, solver settings, \cmd{:WriteGWBudgetFile}, \dots); for an existing model's budget file edit its OC file\\
\cmd{:GridType FILE} needs \cmd{:GridFile} (and the reverse) & give both, or neither\\
\cmd{:GridRotation} applies only to REGULAR and QUADTREE grids & remove it, or change the grid type\\
\cmd{:GridRefinement} size must be positive and smaller than \cmd{:GridCellSize}; RIVERS needs \cmd{:RiverGeometry}; WELLS needs wells & correct the refinement\\
:GridRefinement LAKES found no reservoir whose lake HRU has a polygon & give the reservoir an \cmd{:HRUID} whose polygon is in \cmd{:HRUGeometry}\\
grid cell \dots\ must be a single polygon without holes & fix the cell file\\
GeoJSON file \dots\ is in EPSG:\emph{n} (its "crs" member, a built-in projection) & give \cmd{:InputCRS EPSG:}\emph{n}, or reproject the file to lon/lat\\
GeoJSON file \dots\ has coordinates outside the longitude/latitude range and is in EPSG:\emph{n} \dots, a projection that is not built in & give \cmd{:InputCRS} with the projection's parameters (\cmd{TMERC}, \cmd{ALBERS} or \cmd{LCC}), or reproject the file to lon/lat\\
GeoJSON file / raster \dots\ looks like longitude/latitude: its coordinates \dots\ lie over the HRUs' LATITUDE/LONGITUDE \dots & with \cmd{:InputCRS} give every GeoJSON and raster input in that projection (or remove \cmd{:InputCRS} if every input is in lon/lat)\\
GeoJSON file \dots\ is in EPSG:\emph{n} \dots, but :InputCRS is \dots / is in longitude/latitude, but :InputCRS says the input files are projected & give every input file in the projection of \cmd{:InputCRS}\\
GeoJSON file \dots\ has coordinates outside the longitude/latitude range / raster \dots\ lies outside the longitude/latitude range & projected input: give \cmd{:InputCRS}\\
well \dots\ is given by X and Y & give \cmd{:InputCRS}, or \cmd{LATITUDE} and \cmd{LONGITUDE}\\
layer \dots\ starts dry (warning) & start from \cmd{:GWInitialization STEADY\_STATE}, a smaller \cmd{INITIAL\_HEAD\_DEPTH} or \cmd{:InitialGWHeads} (Section~\ref{ch:layers})\\
HRU mesh exceeds 5 million cells per layer & increase \cmd{:GridCellSize} or the refinement sizes\\
grid exceeds 5 million cells per layer & increase \cmd{:GridCellSize}\\
cannot load MODFLOW 6 library & run \cmd{python tools/get\_mf6.py}, or set \cmd{:MF6Library} or \cmd{RAVEN\_MF6\_LIB}\\
cannot enter MODFLOW 6 folder & output folder removed or not writable during the run\\
\cmd{:SoilZone PHYSICAL}: SOIL[$m$] has zero thickness / no drainable porosity & every coupled HRU with soil needs all soil-zone layers, with \cmd{FIELD\_CAPACITY} $<1$ and \cmd{POROSITY} $>0$\\
soil class \dots\ of SOIL[$M-1$] needs HYDRAUL\_COND & give \cmd{HYDRAUL\_COND} for the class of the bottom soil-zone layer\\
recharge must leave the soil zone from its bottom layer & take percolation to \cmd{GROUNDWATER} from \cmd{SOIL[$M-1$]}\\
RECHARGE\_DELAY \dots\ cannot be used with :SoilZone PHYSICAL & set \cmd{RECHARGE\_DELAY} to 0 in the profiles of soil-zone HRUs; for a delay through a deep unsaturated zone use \cmd{:SoilZoneVadose LAG} (Section~\ref{sec:szvadoselag})\\
the soil zone \dots\ is \dots\ m thick & list physical soil layers only (at most 20 m)\\
Raven is running more time steps than the MODFLOW 6 simulation holds & check \cmd{:Duration} and \cmd{:TimeStep}\\\bottomrule
\end{longtable}

\section{Warnings and stops during a run}\label{app:warnings}
Per-step warnings name the step by its end (``in the step ending \emph{date}'', with the time for
sub-daily steps), the first ten of each kind.
\begin{itemize}
\item ``soil-zone interface did not converge in the step ending \emph{D} (\ldots; continuing with MODFLOW's flows)'', with the
soil-level change and its link, or the failing supply or connection test.
\item ``MODFLOW 6 heads are not finite in the step ending \emph{D} (the solve diverged): \ldots'' (stop).
\item ``MODFLOW 6 head \emph{h} m in cell (\emph{l r c}) in the step ending \emph{D} lies far outside the model's elevations
(\ldots) and the step did not converge: the solve diverged \ldots'' or ``\ldots with MODFLOW's residual \emph{r} m$^3$/d,
\emph{p} \% of the gross flow \emph{G} m$^3$/d (more than 1 \%): its heads are wrong \ldots'' (stops); ``\ldots in a step
MODFLOW did not converge, but its last outer iteration moved the head only \ldots'' and ``\ldots in a converged step'' (warnings).
\item ``\ldots m$^3$ of the negative recharge the top cells could not supply was not held by the soil zone \ldots''.
\item ``\ldots m$^3$ of Raven's recharge was not applied by MODFLOW (a model without NEWTON) on \emph{n} column(s) whose cells are
all inactive (dry) and \emph{m} column(s) whose highest active cell is a specified head \ldots''.
\item ``\ldots m$^3$ of reservoir seepage was not delivered by MODFLOW (cells under the lake that are specified heads or
inactive); returned to the reservoir's subbasin \ldots'' (an existing model).
\item ``boundary \emph{name} has no cells: its head (\emph{h} m, \ldots) lies below the bottom of every cell it selects (highest
\emph{z} m) \ldots'' (at set-up: give a head above those cells, or a \cmd{:BoundaryHead} series that reaches above them).
\item ``:SoilZone PHYSICAL needs an existing model that uses the Newton-Raphson formulation \ldots'' (stop).
\item ``MODFLOW 6 stopped the run during \emph{call}: \emph{the error report of mfsim.lst}'' (stop: an input error MODFLOW finds
inside its library, e.g. \cmd{LINEAR\_ACCELERATION CG} with \cmd{NEWTON}).
\item ``the existing model has no specific yield (no STO package or SY); 0.2 is assumed for the soil-zone checks, the release
limits and the unsaturated columns'' (or ``SPEC\_YIELD \emph{s} of :SoilZoneAquiferClass \emph{class} is used \ldots'' when that
class gives \cmd{SPEC\_YIELD}); ``the existing model's specific yield (STO SY) is below 0.01 in \emph{n} of its \emph{m}
active cells \ldots''.
\item ``:SoilZone PHYSICAL with :SolverMaxOuterIterations \emph{n}'' / ``with OUTER\_MAXIMUM \emph{n} in the model's IMS file
(MODFLOW's MXITER): \ldots a limit below 20 leaves many steps unconverged''.
\end{itemize}

\section{Parameter index}\label{app:paramindex}
Every parameter, its class, and where it is used.
\begin{longtable}{P{0.348\linewidth}P{0.22\linewidth}P{0.36\linewidth}}
\caption{Parameter index: every parameter, its class and where it is used.}\\
\toprule\hdr
Parameter & Class & Used in\\\midrule\endfirsthead
\caption[]{Parameter index: every parameter, its class and where it is used (continued).}\\
\toprule\hdr
Parameter & Class & Used in\\\midrule\endhead
\cmd{K\_HORIZ} & aquifer class & flow between cells (Sec.~\ref{ch:concepts}); well split (Sec.~\ref{ch:wells})\\
\cmd{K\_VERT} & aquifer class & vertical flow between layers (Sec.~\ref{ch:layers})\\
\cmd{SPEC\_STORAGE}, \cmd{SPEC\_YIELD} & aquifer class & storage (Sec.~\ref{ch:concepts}, \ref{ch:layers})\\
\cmd{POROSITY}, \cmd{CLAPP\_B}, \cmd{AIR\_ENTRY\_PRESSURE} & aquifer class & unsaturated columns (Sec.~\ref{sec:szvadose}); \cmd{CLAPP\_B} and \cmd{AIR\_ENTRY\_PRESSURE} also the gap below the soil base (Sec.~\ref{sec:szcaprise})\\
layers, thickness, type & aquifer profile & layer construction (Sec.~\ref{ch:layers})\\
\cmd{INITIAL\_HEAD\_DEPTH} & aquifer profile & initial heads (Sec.~\ref{ch:layers})\\
\cmd{DEFAULT\_BEDROCK\_DEPTH} & aquifer profile & bedrock without a map (Sec.~\ref{ch:layers})\\
\cmd{SOIL\_ZONE\_DEPTH} & aquifer profile & model top (Sec.~\ref{ch:layers}); seepage (Sec.~\ref{ch:seepage})\\
\cmd{SEEPAGE\_LEAKANCE}, \cmd{SEEPAGE\_SMOOTHING\_DEPTH} & aquifer profile & seepage (Sec.~\ref{ch:seepage})\\
\cmd{RIVERBED\_K}, \cmd{RIVERBED\_THICKNESS} & aquifer profile & river conductance (Sec.~\ref{ch:rivers})\\
\cmd{EXTINCTION\_DEPTH} & aquifer profile & water-table ET (Sec.~\ref{ch:et})\\
\cmd{RECHARGE\_DELAY} & aquifer profile & recharge delay (Sec.~\ref{ch:recharge})\\
\cmd{AQUIFER\_PROFILE} & HRU & coupling of HRUs (Sec.~\ref{ch:grid})\\
\cmd{:SeepageParameters} & reservoir & reservoir seepage (Sec.~\ref{ch:reservoirs})\\
\cmd{LEAKANCE} & \cmd{:SurfaceWater\allowbreak Exchange} & wetland and lake exchange (Sec.~\ref{ch:stores})\\
\cmd{CONDUCTANCE}, \cmd{HEAD} & boundary & GHB and CHD (Sec.~\ref{ch:wells})\\
\cmd{SCREEN\_TOP}, \cmd{SCREEN\_BOT}, \cmd{:WellFlowRate} & well & pumping (Sec.~\ref{ch:wells})\\
\cmd{SCREEN\_ELEV}, \cmd{GW\_HEAD} & monitoring well & heads and diagnostics (Sec.~\ref{ch:wells})\\
\cmd{:GridCellSize}, \cmd{:MinCellCoverage} & groundwater option & grid and active cells (Sec.~\ref{ch:grid})\\
\cmd{:GridType}, \cmd{:GridRotation}, \cmd{:GridRefinement}, \cmd{:GridFile}, \cmd{:QuadtreeMaxLevel} & groundwater option & grid types (Sec.~\ref{sec:gridtypes})\\
\cmd{:FlowCorrection}, \cmd{:LayerConnection} & groundwater option & XT3D; layer connections (Sec.~\ref{sec:gridtypes})\\
\cmd{:MF6Simulation}, \cmd{:MF6CRS}, \cmd{:MF6GridFile}, \cmd{:OverlapWeights}, \cmd{:RavenTakesOver}, \cmd{:KeepPackages}, \cmd{:StreamPackage}, \cmd{:SeepageToRaven}, \cmd{:MF6StartDate}, \cmd{:MF6Model} & existing-model option & Chapter~\ref{ch:existing}\\
\cmd{:RiverBedDepth}, \cmd{:SeepageReturnTo} & groundwater option & rivers, seepage (Sec.~\ref{ch:rivers}, \ref{ch:seepage})\\
\cmd{:SolverHead\allowbreak Tolerance}, \cmd{:SolverMax\allowbreak OuterIterations} & groundwater option & Newton solution (Sec.~\ref{ch:mf6})\\
\cmd{:SoilZone}, \cmd{:SoilZoneIteration}, \cmd{:SoilZoneSubgrid}, \cmd{:SoilZoneOutput}, \cmd{:SoilZoneExchangeTiming}, \cmd{:SoilZoneClasses} & soil-zone option & the soil zone and the water table (Sec.~\ref{ch:soilzone})\\
\cmd{:SoilZoneVadose}, \cmd{:SoilZoneVadoseClasses}, \cmd{:SoilZoneVadoseInitialFlux}, \cmd{:SoilZoneCapillaryRise}, \cmd{:SoilZoneAquiferClass} & soil-zone option & unsaturated columns and capillary rise (Sec.~\ref{sec:szvadose}, \ref{sec:szcaprise})\\
\cmd{HYDRAUL\_COND}, \cmd{POROSITY}, \cmd{FIELD\_CAPACITY}, \cmd{STONE\_FRAC}, \cmd{CLAPP\_B}, \cmd{AIR\_ENTRY\_PRESSURE} & soil class (Raven) & soil-zone capacity, conductance and capillary rise (Sec.~\ref{ch:soilzone}, \ref{sec:szcaprise})\\\bottomrule
\end{longtable}

\section{Notation}\label{app:notation}
\begin{longtable}{P{0.2\linewidth}P{0.6\linewidth}P{0.12\linewidth}}
\caption{Notation.}\\
\toprule\hdr Symbol & Meaning & Unit\\\midrule\endfirsthead
\caption[]{Notation (continued)}\\
\toprule\hdr Symbol & Meaning & Unit\\\midrule\endhead
\multicolumn{3}{p{\dimexpr\linewidth-2\tabcolsep}}{\textit{Indices:} HRU $k$ (also $k'$); cell or column $i$ (neighbour $j$); subbasin $p$; MODFLOW layer $\ell$ (also $\ell'$); aquifer profile $a$; boundary $\beta$; polygon corner $v$; soil-zone DEM point $g$; time step $n$ (superscript); Raven soil layer $m$ (\cmd{SOIL[$m$]}).}\\\midrule
$A_k$ & HRU area (\cmd{AREA}, \cmd{.rvh}); $10^6A_k$ in m$^2$ & km$^2$\\
$A_{k,i}$, $A^{\mathrm{cov}}_i$ & overlap of HRU $k$ and column $i$; area of column $i$ covered by coupled HRUs & m$^2$\\
$b$, $\psi_a$ & Campbell exponent and air-entry pressure (\cmd{CLAPP\_B}, \cmd{AIR\_ENTRY\_PRESSURE}; soil classes in mm, aquifer classes in m) & --, m\\
$|\Omega_i|$ & area of cell $i$ (the polygon $\Omega_i$) & m$^2$\\
$b_\ell$, $b_i$, $b^{\mathrm{sz}}$ & layer thickness; thickness of the top active cell of column $i$; saturated thickness of a soil-zone column & m\\
$C$, $C^{\mathrm{sz}}_{k,i}$ & conductance; soil--aquifer conductance of the soil-zone column of HRU $k$ in cell $i$ & m$^2$/d\\
$d_s$, $d_x$, $d_0$ & seepage smoothing depth; ET extinction depth; initial depth to water & m\\
$d_{\mathrm{sz},a}$, $D_a$ & \cmd{SOIL\_ZONE\_DEPTH}, \cmd{DEFAULT\_BEDROCK\_DEPTH} of profile $a$ & m\\
$d_{\mathrm{sz},k}$, $H_m$ & soil-zone thickness of HRU $k$ (\cmd{:SoilZone PHYSICAL}); thickness of \cmd{SOIL[$m$]} & m\\
$\bar d^{\,\mathrm{st}}_i$ & mean depth of water in a surface store over column $i$ & m\\
$\Delta$, $\dt$, $t_{\mathrm{end}}$ & cell size (\cmd{:GridCellSize}); time step; run duration & m, d, d\\
$\delta_b$ & seed spacing along HRU boundaries (HRU mesh) & m\\
$F$, $\bar F_i$, $w_g$ & saturated fraction of a soil-zone column; mean over cell $i$; contact weight of DEM point $g$ & --\\
$h$, $z^{\mathrm{wt}}$ & head; water table of a column (highest wet cell) & m\\
$h^{\mathrm{sz}}$, $z_{b,g}$ & soil water level of a soil-zone column; soil-zone base at DEM point $g$ & m\\
$h_{\mathrm{gw}}$, $h_{\mathrm{res}}$ & local groundwater level of a reservoir; reservoir stage & m\\
$h^{\mathrm{chd}}_i$ & specified head of cell $i$ & m\\
$K$, $K_h$, $K_v$, $K_b$ & hydraulic conductivity; horizontal; vertical; riverbed & m/d\\
$K_{\mathrm{sat},m}$ & saturated conductivity of \cmd{SOIL[$m$]} (\cmd{HYDRAUL\_COND}) & mm/d\\
$k_{\mathrm{seep}}$ & reservoir seepage constant & m$^3$/s/m\\
$L_p$, $L_{p,i}$ & river length of subbasin $p$ on active cells; in cell $i$ & m\\
$L_{\min}$ & shortest unsaturated column (5~cm) & m\\
$\Lambda$, $\Lambda_s$ & leakance (stores); seepage leakance & 1/d\\
$\mathrm{lat}$, $\mathrm{lon}$ & latitude, longitude & degrees\\
$M$, $N_g$ & number of soil-zone layers; DEM points per cell (\cmd{:SoilZoneSubgrid}, squared) & --\\
$N_v$ & number of corners of a polygon & --\\
$N$, $N_c$ & finite volumes (nodes) of an unsaturated column (\cmd{:SoilZoneVadose}, any method); classes of columns (\cmd{:SoilZoneClasses}) & --\\
$\eta_m$, $S_{\mathrm{fc},m}$, $\mathit{SF}_m$, $S_d$, $\chi_m$ & porosity, field capacity, stone fraction, drainable porosity and saturated share of \cmd{SOIL[$m$]} & --\\
$\phi^{\mathrm{st}}_k$ & water in a Raven surface store of HRU $k$ & mm\\
$\phi_{k,i,m}$, $\phi_{\max,m}$, $\phi_{\mathrm{fc},m}$ & water of \cmd{SOIL[$m$]} in the soil-zone column of HRU $k$ in cell $i$; capacity; field capacity & mm\\
$Q$, $q$, $q^{\mathrm{sz}}_{k,i}$ & flow (positive into the aquifer for boundaries); soil-zone exchange of a column & m$^3$/d\\
$Q^{\mathrm{out}}_p$ & channel flow of subbasin $p$ & m$^3$/s\\
$q_{\mathrm{top}}$, $q_{\mathrm{bot}}$ & flux into the top and out of the bottom of an unsaturated column & m/d\\
$R_\oplus$ & radius of the Earth (sphere) & m\\
$r_k$, $r_i$ & recharge depth of HRU $k$ in a step; recharge rate of cell $i$ & mm, m/d\\
$s$ & river stage & m\\
$S_s$, $S_y$ & specific storage; specific yield & 1/m, --\\
$\mathcal{S}_i$ & storage capacity of cell $i$ & m$^2$\\
$T$ & transmissivity & m$^2$/d\\
$\theta$, $\theta_s$, $\psi$ & water content and saturated content (\cmd{POROSITY}) of an unsaturated column; pressure head (negative: suction) & --, --, m\\
$\tau$ & recharge delay time constant & d\\
$V$ & volume & m$^3$\\
$V^{\mathrm{dl}}_k$, $V^{\mathrm{carry},n}_p$, $V^{\mathrm{st}}_{k,i}$ & recharge delay store; river loss carried out of step $n$; store water of HRU $k$ in cell $i$ & m$^3$\\
$V^{\mathrm{rej}}$, $V^{\mathrm{sz}}_{\mathrm{applied}}$, $V^{\mathrm{short}}$ & rejected recharge; soil-zone exchange applied to Raven; soil-zone shortfall & m$^3$\\
$V^{\mathrm{uz}}_{\mathrm{in}}$, $V^{\mathrm{uz}}_{\mathrm{to\,aq}}$, $V^{\mathrm{uz}}_{\mathrm{cap}}$, $V^{\mathrm{wt}}_{\mathrm{cap}}$, $V^{\mathrm{debt}}$ & unsaturated columns: inflow, outflow to the aquifer, capillary rise into the soil zone; capillary rise from the water table; lagged-exchange debt & m$^3$\\
$z_{\mathrm{land}}$, $z_{\mathrm{top}}$, $z_{\mathrm{br}}$, $z_{\mathrm{bed}}$ & land surface, model top, bedrock, riverbed elevations & m\\\bottomrule
\end{longtable}

